% Chapter 5
\chapter{Implementation Plan and Execution Strategy}
\label{chap:implementation}

This chapter describes how the proposed observability architecture will be implemented in a controlled and reproducible manner. The focus is on deployment strategy, execution sequence, and artifact traceability, ensuring that implementation activities remain aligned with requirements and design decisions.

\section{Implementation Objectives}

The implementation phase operationalizes the selected architecture by deploying the telemetry pipeline, configuring data collection, and validating end-to-end data flow across services. The phase prioritizes reproducibility, incremental rollout, and risk-aware execution.

\section{Implementation Work Packages}

Implementation is organized into progressive work packages:
\begin{enumerate}
    \item \textbf{Environment preparation:} baseline configuration validation, network checks, and service inventory.
    \item \textbf{Observability stack deployment:} deployment of selected components for metrics, logs, traces, and visualization.
    \item \textbf{Instrumentation configuration:} service-level telemetry configuration and correlation metadata setup.
    \item \textbf{Operational hardening:} access control, retention policies, and backup/recovery safeguards.
    \item \textbf{Reproducibility assets:} scripts, templates, and runbooks for repeatable execution.
\end{enumerate}

\section{Traceability and Quality Gates}

Each implementation increment is validated by explicit quality gates:
\begin{itemize}
    \item successful deployment from documented procedures;
    \item telemetry visibility for each target service;
    \item basic cross-signal correlation capability;
    \item documented rollback and recovery procedures.
\end{itemize}

This gate-based approach ensures that implementation outcomes can be audited and replicated during evaluation.

\section{Chapter Summary}

This chapter defined the implementation strategy required to move from architecture to a working observability deployment. The next chapter specifies the evaluation framework, including a-priori metrics and measurement procedures.
